% A283 -- Thermodynamik in der FFGFT: ein Gesamtüberblick
% Johann Pascher -- ORCID 0009-0000-6518-4064
% Kapitelquelle zur A-Serie; Wrapper: wr_standalone_A4/A283_Thermodynamik_FFGFT_De.tex
% Status: August 2026

\providecommand{\qedsq}{\raisebox{0.03em}{\framebox[0.62em]{\rule{0pt}{0.42em}}}}

% ===================== KURZFASSUNG =====================
\begingroup\small\noindent
\textbf{Kurzfassung.}
Die Thermodynamik ist in der FFGFT auf fünf Ebenen verankert ---
nicht als Postulat, sondern als Konsequenz der Geometrie.
Dieses Dokument fasst die verteilten Einzelresultate zusammen,
ordnet sie in eine einheitliche Argumentationskette und benennt,
wo die Ebenen ineinandergreifen und wo sie getrennt bleiben.
Ein Abschlussabschnitt räumt ein, dass beide Rahmenbeschreibungen ---
expandierender Raum und fraktale Verfeinerung --- Modelle sind,
die dieselben Beobachtungen strukturell entgegengesetzt deuten.
\par\endgroup

\bigskip

\begingroup\small\noindent
\textbf{Schichtmarkierungen.}
\textbf{[SETZUNG]} deklarierter Ausgangspunkt ---
\textbf{[B]} algebraisch/mathematisch aus deklarierten Grundlagen abgeleitet ---
\textbf{[K]} numerisch aus Korpus-Konstanten berechnet ---
\textbf{[Q]} korpusexterne Primärquelle ---
\textbf{[S]} Plausibilitätsskizze ---
\textbf{[H]} offene Forschungsfrage.
\par\endgroup

\bigskip

% ===================== 1 =====================
\section{Warum ein Gesamtdokument}

Der thermodynamische Inhalt der FFGFT ist auf mehrere Dokumente
verteilt entstanden: Dok.~157 behandelt den Zeitpfeil, A271--A282
die Landauer-Buchhaltung, Dok.~313 die kosmologische Thermodynamik.
Jedes Dokument ist in sich geschlossen; keines gibt einen Überblick
darüber, wie die Teile zusammenhängen.

\textbf{[B]} Drei Fragen bleiben ohne Gesamtdokument unbeantwortet:
(1) Wo berühren sich die Ebenen, und wo bleiben sie getrennt?
(2) Welche Aussagen hängen am $\xi$-Schema, welche sind davon unabhängig?
(3) Was ist der epistemische Status jeder Aussage?
Dieses Dokument beantwortet diese Fragen und enthält keine
neuen inhaltlichen Behauptungen.

% ===================== 2 =====================
\section{Ebene 1 --- Der Zeitpfeil ist geometrisch (Dok.~157)}

\subsection{Das Standardproblem}

Alle fundamentalen Gleichungen der Physik sind zeitumkehrinvariant.
Der zweite Hauptsatz ($\Delta S \ge 0$) ist statistischer Natur
und verlagert das Problem auf die unerklärte niedrige Anfangsentropie
(Penrose: $\sim 10^{-10^{123}}$).

\subsection{Die FFGFT-Antwort}

\textbf{[B]} Drei geometrische Argumente greifen ineinander (Dok.~157):

\textbf{Argument 1: Algebraische Asymmetrie.}
$\tilde{T} \cdot m = 1$ erzwingt $\tilde{T} > 0$, weil Masse definit
positiv ist. Die Zeitumkehr würde $m \to -m$ verlangen ---
physikalisch unmöglich.

\textbf{Argument 2: Topologische Wicklungsrichtung.}
Das Vakuumfeld wickelt sich mit $\dot\theta = m = 1/\tilde{T}$
entlang des $T^4$. Diese Chiralität ist topologisch geschützt.

\textbf{Argument 3: Fraktale Dimensionsasymmetrie.}
$D_f = 3 - \xi_0 \approx 2{,}9999$ bricht die exakte Spiegelsymmetrie;
Pfadintegrale für vorwärts und rückwärts unterscheiden sich um
$O(\xi_0)$.

\subsection{Konsequenz}

\textbf{[B]} Die Logik kehrt sich um: $\Delta S \ge 0$ ist in der FFGFT
eine \emph{Konsequenz} der Geometrie, kein Postulat. Die
\enquote{Past Hypothesis} wird unnötig --- die niedrige Anfangsentropie
ist geometrisch erzwungen, keine Anfangsbedingung.

% ===================== 3 =====================
\section{Ebene 2 --- Landauer als Phasenraum-Buchhaltung (A271)}

\subsection{Die Phasenraum-Formulierung}

\textbf{[SETZUNG]} Ein physikalisches System belegt ein Gebiet im
Phasenraum der Größe $W$ (gemessen in $h^f$). Das Gebiet existiert
unabhängig davon, wie man es benennt.

\textbf{[B]} Eine Abbildung, die mehrere Ausgangsgebiete $\{A, B\}$
auf ein Zielgebiet $Z$ reduziert, verkleinert das Volumen:
$\Delta S_\text{System} = -k_B \ln(W_{A+B}/W_Z) \le 0$.
Der Satz von Liouville verbietet Volumenschrumpfung im abgeschlossenen
System. Ein Wärmebad muss das Volumen aufnehmen:
\[
  \boxed{Q \;\ge\; k_B\,T\,\ln\!\left(\frac{W_\text{vorher}}{W_\text{nachher}}\right)}
\]
mit $\ln 2$ als Binär-Spezialfall.

\subsection{Drei Trennungen}

\textbf{[B]}
(1) \emph{Bit $\ne$ physikalisches Gebiet.} Shannon- und Boltzmann-Entropie
sind nur bei Gleichverteilung über physikalische Mikrozustände identisch.
(2) \emph{Inhaltsneutralität.} Die Thermodynamik zählt Gebiete, liest sie
nicht. Bit \enquote{wahr} und Bit \enquote{falsch} kosten identisch.
(3) \emph{Drei Ebenen getrennt.} Konstruktion / Zustandsübergang / Beschreibung.
Landauer-Kosten fallen auf Ebene~1 und~2 an, nicht auf~3.

% ===================== 4 =====================
\section{Ebene 3 --- Phasenraum-Buchhaltung universell (A280--A282)}

\textbf{[B]} A280: Das Gebiet kommt zuerst. Ein Glas ins Weltmeer ---
das Volumen verschwindet nicht, verteilt sich auf
$1{,}4 \times 10^9$\,km³. Der Preis: das Bad muss das Volumen
aufnehmen, Wärme fließt. Zählbasis ist Unterscheidbarkeit
(Orthogonalität), nicht Energie; absolute Grunddiskretisierung: $h^f$.

\textbf{[B]} A281: Landauers eigene Einschränkung --- logisch irreversibel
ist eine Beschreibungseigenschaft, nicht eine physikalische.
Hinreichend für Wärme ist allein die Nicht-Bijektivität der
Gebietsabbildung. Bennetts Resultat: Jede logisch irreversible Funktion
ist als Folge bijektiver Abbildungen ausführbar; Kosten kleben am
\emph{Vergessen}, nicht an der Operation.

\textbf{[B]} A282: Der Token-Rechner als Grenzfall. Buchhaltung gilt exakt,
Landauer-Schranke gilt dort \emph{nicht} (Ingenieurskosten dominieren).
Grenzübergang zur Brownschen Rechenkugel verbindet stetig mit Landauers
Geltungsbereich.

\textbf{[H]} Offene Frage: Minimaler Kohärenzpreis für Quantenkohärenz
unabhängig von der Fehlerkorrektur (A271, §\,7.4).

% ===================== 5 =====================
\section{Ebene 4 --- Kosmologische Thermodynamik (Dok.~312/313)}

\subsection{Thermische Geschichte als geometrische Bedingung}

\textbf{[K]} Dok.~313: Drei Phasen auf dem beobachtbaren Halbzyklus
$(\tau_c/2 \approx 46\,\text{Gyr})$:

\begin{center}
\small
\begin{tabular}{llll}
\toprule
Phase & Bereich & $\theta/\pi$ & Charakter \\
\midrule
Strahlungsdomination & $z > 3000$ & $> 0{,}90$ &
  $m \to m_\text{min}$, Uhren langsamst \\
Übergang & $z \approx 3000$--$1090$ & $0{,}90$--$0{,}979$ &
  Rekombination, Streuwand \\
Materie-Dominanz & $z < 3000$ & $< 0{,}90$ &
  $m(\theta)$ wächst, Strukturbildung \\
\bottomrule
\end{tabular}
\end{center}

\subsection{Die fünfstufige Kette}

\textbf{[K]} Vollständige Herleitung von $\Omega_m^*$ aus $\xi_0$ allein:
\begin{center}
\small
\begin{tabular}{clc}
\toprule
& Schritt & Ergebnis \\
\midrule
(1) & $H_0 = (\pi/2)\,c\,\xi_0^{10}/\bar\lambda_e$ & $h = 0{,}6682$ \\
(2) & $k_BT_0 = (16/9)\,\xi_0$ (Dok.~061) & $T_0 = 2{,}751$\,K \\
(3) & $\Omega_r \propto T_0^4/H_0^2$ & $9{,}643\times10^{-5}$ \\
(4) & $K_\text{frak} = 1-100\,\xi_0$ (A040) & $0{,}98667$ \\
(5) & $\int_0^\infty dz/E(z) = \pi/K_\text{frak}$ &
  $\boxed{\Omega_m^* = 0{,}3136}$ \\
\bottomrule
\end{tabular}
\end{center}
Planck: $\Omega_m = 0{,}315 \pm 0{,}007$ ($-0{,}19\,\sigma$).
Kein einziger Schritt enthält einen an $\Omega_m$ angepassten Parameter.
Die Strahlung ($\Omega_r \propto T_0^4/H_0^2$) ist das Bindeglied
zwischen Teilchen- und kosmischem Sektor.

\subsection{Überbestimmung und offener Kern}

\textbf{[K]} $T_0$ ist über zwei unabhängige Wege bestimmt:
Teilchensektor (Weg~1: $k_BT_0 = (16/9)\xi_0$) und kosmischer Sektor
(Weg~2: Antipodenbedingung). Beide liefern $\Omega_m$ auf $0{,}1\,\%$
übereinstimmend --- dasselbe Überbestimmungsmuster wie $\alpha$ in A130.

\textbf{[K]} Offen (D(ii)): Grobkörnige Entropie auf der
Rückkehr-Hälfte. Poincaré-Diskrepanz: $\sim 10^{10^{104}}$ dynamische
Zeiten gegen $\tau_c \approx 92$\,Gyr. Periodizität ist eine
Randbedingung, keine Folge der Dynamik. Drei Fälle (A050/Dok.~295):
Fall~A ($d=0$, 1 Umlauf), Fall~B ($\xi$ eingefroren, 75 Umläufe),
Fall~C ($\xi$ läuft, Logspirale --- hebt Entropieforderung auf,
kostet aber $\tau_c$).

\subsection{Falsifikationspunkte}

\textbf{[K]} Low-$\ell$-Cutoff: $\ell_\text{min} = 2\pi D_\text{LSS}/L^* = 3{,}08$
--- trifft die seit COBE (1992) unerklärte Quadrupol-Unterdrückung.
Hauptfalsifikationstest: $m_\tau = 1776{,}97$\,MeV (Belle-II).

% ===================== 6 =====================
\section{Ebene 5 --- Was die Ebenen trennt}

\textbf{[B]} Die Landauer-Buchhaltung (Ebenen~2 und~3) berührt das
$\xi$-Schema nicht. Sie gilt in jedem System mit Hamiltonscher oder
unitärer Dynamik --- unabhängig von der Torusgeometrie. Explizit
gebucht in A271: \enquote{Berührung mit dem $\xi$-Schema der FFGFT:
keine.}

Der Zeitpfeil (Ebene~1) dagegen ist direkt an $D_f = 3-\xi_0$ und
$\tilde{T} \cdot m = 1$ gebunden --- FFGFT-spezifisch.
Die kosmologische Thermodynamik (Ebene~4) hängt vollständig an
$\xi_0$, $H_0$ und der Torusgeometrie.

\begin{center}
\small
\begin{tabular}{@{}p{0.30\textwidth}p{0.20\textwidth}p{0.22\textwidth}p{0.18\textwidth}@{}}
\toprule
\textbf{Ebene} & \textbf{Dokument} & \textbf{$\xi$-Abhängigkeit} & \textbf{Status} \\
\midrule
Zeitpfeil & Dok.~157 & direkt & [B]/[S] \\
Landauer & A271 & keine & [B] \\
Buchhaltung universell & A280--A282 & keine & [B] \\
Kosmol.\ Thermodynamik & Dok.~312/313 & direkt & [K] \\
Informationsphysik & Dok.~243--301 & indirekt & [S]/[B] \\
\bottomrule
\end{tabular}
\end{center}

% ===================== 7 =====================
\section{Argumentationskette: von der Geometrie zur Wärme}

\textbf{[B]} Die fünf Ebenen bilden eine Kette von der FFGFT-Geometrie
zur klassischen Thermodynamik:
(1) $\xi_0$ und $T^4$ legen die Geometrie fest.
(2) $\tilde{T} \cdot m = 1$ und die Toruswicklung definieren
geometrisch eine Zeitrichtung.
(3) In dieser gerichteten Zeit gibt es mehr Phasenraumvolumen in
Vorwärtsrichtung --- $\Delta S \ge 0$ als geometrische Konsequenz.
(4) Nicht-bijektive Phasenraumabbildungen (Landauer) kosten Wärme ---
unabhängig von $\xi_0$, im gleichen Phasenraum-Bild.
(5) Die thermische Geschichte ($\Omega_m^*$, $\Lambda^*$) folgt aus
$\xi_0$ direkt.

Der Satz von Liouville und der zweite Hauptsatz sind in dieser
Kette \emph{Konsequenzen}, keine Postulate.

% ===================== 8 =====================
\section{Informationsthermodynamik im IPI-Kontext (Dok.~243--301)}

\textbf{[S]} Eine weitere Gruppe von Dokumenten verbindet die
FFGFT-Thermodynamik mit der Informationsphysik-Community (IPI-Kontext):
Dok.~251 (Vergleich mit Vopsons Informationsmassen-Hypothese),
Dok.~247 (Landauers Prinzip im Kontext von Informationszuständen),
Dok.~290 (Informationsfrage Betrag/Phase),
Dok.~301 (Informations-Umkehrtest).
Diese Dokumente sind Plausibilitätsskizzen oder teilweise algebraische
Herleitungen; sie erschließen den Anschluss an aktuelle Debatten,
ohne neue Kernresultate zur FFGFT-Thermodynamik beizutragen.

% ===================== 9 =====================
\section{Phasen, Strahlung und Temperatur im Zeitzyklus (Dok.~313 im Detail)}

\subsection{Gradient auf einem Kreis}

\textbf{[B]} Die beobachtbare Geschichte zeigt eine klare thermische
Abfolge: Strahlungsdomination, Übergang, Materie-Dominanz. Das
erweckt den Eindruck eines gerichteten Prozesses mit einem Anfang.
Die FFGFT-Antwort (Dok.~313): Der Gradient ist keine Entwicklung
\emph{aus} einem Anfang, sondern eine Landkarte \emph{entlang} des
Zeitzyklus --- wie Klimazonen entlang eines Meridiankreises.

\subsection{Drei Phasen auf dem Halbzyklus}

\textbf{[K]} Die beobachtbare Geschichte füllt einen Halbzyklus
($\tau_c/2 \approx 46$\,Gyr). Zykluspositionen:

\begin{center}
\small
\begin{tabular}{lccc}
\toprule
& & \multicolumn{2}{c}{$\theta/\pi$} \\
\cmidrule(l){3-4}
Ereignis & $z$ & glatt & fraktal \\
\midrule
Materie-Strahlung-Gleichheit & $\sim\!3400$ & $0{,}93$ & $0{,}92$ \\
Streuwand (Rekombination) & $1090$ & $0{,}992$ & $\mathbf{0{,}979}$ \\
Partikelhorizont & $z_\text{max}$ & $1{,}012$ & $\mathbf{0{,}9986}$ \\
\bottomrule
\end{tabular}
\end{center}

\subsubsection{Phase~1 --- Strahlungsdomination (nahe dem Antipoden)}

In der Massenlauf-Lesart: Massen $m(\theta)$ sind minimal,
intrinsische Zeiten $\tilde{T} = 1/m$ maximal --- Uhren laufen am
langsamsten. Die Endpunktsteigung des Massenprofils am Antipoden:
\[
  \left|\frac{dm}{d\theta}\right|_\text{Antipode}
  = \sqrt{\Omega_r} \approx 0{,}0096
\]
Kein freier Parameter --- die beobachtete Strahlungsdichte erscheint
hier als \emph{Glattheitsbedingung} der periodischen Schließung.

\subsubsection{Phase~2 --- Materie-Strahlung-Übergang (z $\approx$ 3000 bis 1090)}

Mit der fraktalen Wegkorrektur $K_\text{frak} = 1-100\xi_0$ (A040,
unabhängig abgeleitet) sitzt der Partikelhorizont auf $0{,}14\,\%$
am Antipoden --- Verbesserung um Faktor~9 gegenüber der glatten
Rechnung ($1{,}2\,\%$), ohne angepassten Parameter.

Die Wegkorrektur greift bei Weg/Struktur, kürzt bei Weg/Weg.
Invarianztest: CMB-Peakpositionen bleiben unverändert
(Test bestanden ohne dafür konstruiert worden zu sein).

\subsubsection{Phase~3 --- Materie-Dominanz (heutige Epoche)}

Energiedichte $\propto (1+z)^3$. In der Massenlauf-Lesart wachsen
die Massen mit sinkendem $z$; Uhren laufen schneller; Strukturbildung
entfaltet sich.

\subsection{Die Temperatur als Strukturgröße}

\textbf{[K]} Im Standardmodell: $T(z) = T_0(1+z)$ --- adiabatische
Rückextrapolation. In der FFGFT: $T_0$ stammt aus der Atomskala
(Dok.~061):
\[
  k_B T_0 = \tfrac{16}{9}\,\xi_0 \approx 2{,}370\times10^{-4}\,\text{eV}
  \quad(\text{Messung: }2{,}349\times10^{-4}\,\text{eV, }{+0{,}92\,\%})
\]
Die CMB-Temperatur ist eine \emph{Strukturgröße} des $\xi$-Schemas,
kein Endpunkt einer Abkühlungsgeschichte.

\textbf{[K]} Reichweitengrenze (Dok.~313): Gleichgewichte an anderen
Orten des Zeitzyklus sind aus dem heutigen \emph{nicht} ableitbar
durch Fortschreiben. Unberührt: die $\Omega_m^*$-Kette (läuft über
$\xi_0$, nicht über adiabatische Rückrechnung).

\subsection{Überbestimmung von $T_0$}

\textbf{[K]} Zwei unabhängige Wege:
Weg~1 (Teilchensektor): $k_BT_0 = (16/9)\xi_0$, kein kosmologischer Eingang.
Weg~2 (kosmischer Sektor): Antipodenbedingung, kein teilchenphysikalischer Eingang.
Beide liefern $\Omega_m$ auf $0{,}1\,\%$ übereinstimmend
($0{,}3136$ vs.\ $0{,}3139$) --- dasselbe Muster wie $\alpha$ in A130.
Hebel: $dT_0/T_0 \big/ d\Omega_m/\Omega_m \approx -12$.

\subsection{Untergrenze und Hawking}

\textbf{[K]} Energiequant des Zeitzyklus (Dok.~312):
$E_\text{min} = \hbar H_0 = 2{,}28\times10^{-52}$\,J.
Endliches $z_\text{max} = m_e/m_\text{min} - 1 = 3{,}59\times10^{38}$.
Singularität ausgeschlossen: Massen können nicht gegen null laufen.

\textbf{[K]} Eine KMS-Regel verbindet Gibbons--Hawking, Unruh und
Hawking: $T = \hbar/(k_B\tau)$ mit $\tau$ als lokaler Zeitperiode.
Die kompakte Zeitrichtung liefert $T_\text{GH}$ ohne Horizont.
Familienleiter trifft Kosmos bei $r_s = R_H/2$ (exakt $0{,}500$).
Hawking-Verdampfung für $M_\odot$: $56$ Größenordnungen zu langsam
für $\tau_c$ --- verschärft das Entropie-Hindernis D(ii).

\subsection{Low-$\ell$-Cutoff}

\textbf{[K]} Auf einem Torus der Kantenlänge $L^*$ fehlen Moden mit
Wellenlänge $> L^*$:
\[
  \ell_\text{min} = \frac{2\pi D_\text{LSS}}{L^*} = 3{,}08
\]
Kein angepasster Parameter. Trifft die Quadrupol-Unterdrückung
($\ell = 2$, Faktor~$\sim 1/6$), seit COBE (1992) unerklärter Befund
in $\Lambda$CDM. Ordnungsaussage [K$|$P20]; genaue Modensummation
auf dem Torus steht noch aus.

\subsection{Belastbare und nicht-belastbare Aussagen}

\textbf{[K]} Drei Klassen (Dok.~313):

\begin{center}
\small
\begin{tabular}{p{3.2cm}p{4cm}p{3cm}}
\toprule
Klasse & Beispiele & Belastbarkeit \\
\midrule
Zählungen & $|\text{Aut}(D4)| = 1152$, Kusszahl~24 & exakt \\
Verhältnisse gleicher Ebene & $\alpha$, $m_\mu/m_e$, Koide~$Q$ &
  bis $10^{-10}$ \\
Mit Skalenbrücke & $G$, $H_0$, absolute Massen & nur mit Anker \\
Zahl- und Dauergrößen & $\eta = n_B/n_\gamma$, Fensterdauern &
  nicht übertragbar \\
\bottomrule
\end{tabular}
\end{center}

BBN-Observable ($Y_p$, D/H, ${}^7$Li): kein reines Verhältnis
gleicher Ebene --- dasselbe gilt für $\Lambda$CDM.

% ===================== 10 =====================
\section{Massebildung, fraktale Verfeinerung und der Schein der Ausdehnung}

\subsection{Was am Antipoden fehlt: metrische Schärfe}

\textbf{[K]} Alle physikalischen Längen die an Masse hängen, skalieren
mit $1/m$: Compton-Wellenlänge $\bar\lambda = \hbar/(mc)$,
Bohr-Radius $a_0 \propto 1/m_e$, intrinsische Zeit $\tilde{T} = 1/m$.
Am Antipoden ($m \to m_\text{min}$) laufen alle gegen ihr Maximum.
Ohne lokalisierbare massebehaftete Objekte fehlt \emph{metrische Schärfe}.

\textbf{[H]} Was metrische Struktur bei $m \to m_\text{min}$ präzise
bedeutet, ist im Korpus noch nicht ausgearbeitet.

\subsection{Photonen: kein Einfrieren}

\textbf{[B]} Photonen tragen keine Ruhemasse --- $\tilde{T} \cdot m = 1$
gilt für sie nicht direkt. Ein Photon am Antipoden friert nicht ein;
es propagiert mit $c$. Die CMB-Photonen verlieren Energie geometrisch
auf dem Weg zu uns (A265):
\[
  \frac{dE}{dx} = -\xi\,E
  \;\Rightarrow\;
  1 + z = e^{\xi x}
\]
Achromatisch: alle Wellenlängen gleich gestreckt.
Aus unserer Sicht läuft alles gleichmäßig. Der Blick zurück ist ein
Blick auf einen Ort des Zeitzyklus, kein Blick in eine Vergangenheit.

\subsection{Ausdehnung als fraktale Verfeinerung nach innen}

\textbf{[B]} Im Standardmodell: Raum expandiert nach außen,
Skalenfaktor $a(t)$ wächst. In der FFGFT gibt es keinen expandierenden
Raum. Mit zunehmendem $m(\theta)$ werden alle charakteristischen
Längen ($\propto 1/m$) kleiner. Materie wird strukturierter, feiner
aufgelöst. Das ist \textbf{Verfeinerung nach innen}, kein Wachstum
nach außen.

\textbf{[B]} $K_\text{frak} = 1-100\xi_0$ (A040): Dimensionsdefizit
$D_f = 3-\xi_0$, lokal $\sim 0{,}003\,\%$, kumulativ über 17
Größenordnungen $\sim 100$.

\textbf{[B]} Rekursion $\xi_{n+1} = \xi_n(1-100\xi_n)$ (A050),
Fixpunkt: null. Aufgerollt: Windung schließt nach 75 Umläufen
($\tau_c$-Zyklus). Entrollt: logarithmische Spirale
$\rho = 75\,e^{\beta/2\pi}$ --- was das Standardmodell Expansion nennt.

\textbf{[K]} Beide Beschreibungen sind beobachtungsäquivalent (A265).
Ruler und Uhren ändern sich im gleichen Verhältnis; kein lokaler
Beobachter merkt die Skalierung. Was wir als Expansion sehen, ist der
Unterschied zwischen unseren heutigen Massen und dem rotverschobenen
Licht vom Antipoden.

\textbf{[S]} Bild: Mandelbrot-Zoom. Zoom nach innen enthüllt immer
feinere Strukturen, kein Wachstum nach außen. Antipode = unzoomte
Ansicht (minimale Differenzierung, reine Strahlung). Heutige Position
= hochgezoomte Ansicht (scharfe Atom- und Sternstrukturen).
Das Standardmodell dreht das Bild um --- eine Koordinatenwahl,
keine physikalische Aussage.

% ===================== 11 =====================
\section{Offene Punkte}

\textbf{[H]} Drei Fragen bleiben offen:
(1) Minimaler Kohärenzpreis für Quantenkohärenz unabhängig von der
Fehlerkorrektur (A271, §\,7.4).
(2) Kontinuierlicher Entropiestrom als Funktion von Takt, Bitbreite
und logischer Irreversibilität (A271, §\,8).
(3) Metrische Struktur im Grenzfall $m \to m_\text{min}$ --- ob
ein Regime reiner Feldgeometrie ohne lokalisierbare Teilchenmoden
existiert.

% ===================== 12 =====================
\section{Schichtstatus-Zusammenfassung}

\begin{center}
\small
\begin{longtable}{@{}p{0.70\textwidth}l@{}}
\toprule
\textbf{Aussage} & \textbf{Status} \\
\midrule
\endfirsthead
\toprule
\textbf{Aussage} & \textbf{Status} \\
\midrule
\endhead
$\tilde{T} > 0$ erzwingt geometrischen Zeitpfeil & \textbf{[B]} \\
Toruswicklung: topologisch geschützte Zeitrichtung & \textbf{[B]} \\
$D_f = 3-\xi_0$ bricht Zeitumkehrsymmetrie & \textbf{[B]} \\
$\Delta S \ge 0$ als geometrische Konsequenz & \textbf{[B]} \\
Landauer: $Q \ge k_BT\ln 2$ aus Liouville & \textbf{[B]} \\
Inhaltsneutralität der Thermodynamik & \textbf{[B]} \\
Bennetts reversible Logik & \textbf{[B]} \\
$\Omega_m^* = 0{,}3136$ aus $\xi_0$-Kette & \textbf{[K]} \\
$\Lambda^* = 1/R_H^2$ geometrischer Träger & \textbf{[K]} \\
Überbestimmung $T_0$: zwei Sektoren auf $0{,}1\,\%$ & \textbf{[K]} \\
Singularität ausgeschlossen durch $\hbar H_0$ & \textbf{[K]} \\
$\ell_\text{min} = 3{,}08$ (Low-$\ell$-Cutoff) & \textbf{[K]$|$P20} \\
Langsamste Uhren am Antipoden & \textbf{[K]} \\
Compton-Wellenlänge, Bohrradius $\propto 1/m$ & \textbf{[B]} \\
Photonen achromatisch rotverschoben: $1+z = e^{\xi x}$ & \textbf{[B]} \\
Ausdehnung = fraktale Verfeinerung nach innen & \textbf{[B]} \\
Rekursion: Fixpunkt null, Spirale im entrollten Fall & \textbf{[B]} \\
Beobachtungsäquivalenz mit $\Lambda$CDM & \textbf{[K]} \\
Landauer berührt $\xi$-Schema nicht & \textbf{[B]} \\
Grobkörnige Schließung D(ii): offen & offen \\
$m=0$ strukturell ausgeschlossen: kein leerer Zustand (A060) & \textbf{[B]} \\
Doppelcutoff: $L_0 \Rightarrow E_\text{max}$, $L^* \Rightarrow E_\text{min}$ & \textbf{[K]} \\
$m_\text{min} = \hbar H_0/c^2$ aus IR-Cutoff (Dok.~312/313) & \textbf{[K]} \\
Lokale Feldmoden am dünnsten Punkt auflösbar? & \textbf{[H]} \\
Past Hypothesis als geometrische Notwendigkeit & \textbf{[S]} \\
Mandelbrot-Zoom als Bild & \textbf{[S]} \\
\bottomrule
\end{longtable}
\end{center}

\textbf{Bilanz:} 14$\times$ \textbf{[B]} $\cdot$
8$\times$ \textbf{[K]} $\cdot$
3$\times$ \textbf{[S]} $\cdot$
2$\times$ \textbf{[H]} $\cdot$
1$\times$ offen.

% ===================== 13 =====================
\section{Abschlussbemerkung: Zwei entgegengesetzte Modelle}

\textbf{[S]} Eine ehrliche Einräumung gehört an den Schluss.

Für uns --- als Beobachter mit unseren heutigen Massen und Uhren
--- zählt natürlich unsere Sicht. Wir denken in Lichtwegen. Wir
messen Rotverschiebungen. Wir erleben Zeit als vorwärts laufend.
Das ist real in dem Sinne, dass es das ist, was wir messen, und
unsere Instrumente darauf kalibriert sind.

Das Standardmodell ($\Lambda$CDM) nimmt genau diese Sicht als
Ausgangspunkt: expandierender Raum, Urknall als Anfang,
Rotverschiebung als Raumstreckung, CMB als Schnappschuss einer
frühen Epoche. Es ist ein Modell, das mit den Beobachtungen
hervorragend übereinstimmt.

Die FFGFT schlägt eine andere Koordinatenwahl vor: kein
expandierender Raum, kein Urknall als Ereignis, Rotverschiebung
als geometrischer Energieverlust, CMB als Blick auf einen Ort des
Zeitzyklus. Auch diese Beschreibung stimmt mit den Standardtests
überein (A265, kosmologische Entartung).

\textbf{[B]} Was die beiden Modelle unterscheidet:

\begin{itemize}[leftmargin=1.4em,itemsep=0.2em]
  \item $\Lambda$CDM verallgemeinert von der Beobachterperspektive
        nach \emph{außen}: expandierender Raum.
  \item FFGFT verallgemeinert von der Geometrie nach \emph{innen}:
        kompakter Torus, fraktale Verfeinerung.
\end{itemize}

Beide sind Modelle. Beide verlassen die direkte Beobachtung und
fügen eine Interpretationsschicht hinzu.

\textbf{[S]} Was bleibt: zwei strukturell entgegengesetzte Wege,
dasselbe zu beschreiben. Im einen dehnt sich alles nach außen aus,
im anderen verfeinert sich alles nach innen. Im einen hat die Zeit
einen Anfang, im anderen ist sie ein geschlossener Zyklus. Im einen
ist die CMB ein Fenster in die Vergangenheit, im anderen ein Fenster
auf einen Ort.

Das ist kein Schönheitsfehler --- es ist eine offene Frage, präzise
genug formuliert, um entscheidbar zu sein. Die Falsifikationstests
($m_\tau$, $\ell_\text{min}$, $\Omega_m^*$-Kette) werden zeigen,
ob die geometrische Innenperspektive der FFGFT mehr trägt als die
expandierende Außenperspektive --- oder weniger.

% ===================== QUELLEN =====================
\section{Quellverweise}

\begin{description}[leftmargin=3.2em,style=sameline,itemsep=0.3em]
\small
\item[{Dok.~061}] Pascher, J., \emph{Temperatureinheiten und CMB in der
T0-Theorie}. Altkorpus, 2025.
\item[{Dok.~157}] Pascher, J., \emph{Der Zeitpfeil in der T0-Theorie}.
Altkorpus, 2025.
\item[{Dok.~312}] Pascher, J., \emph{Die Abschlussskala}. August 2026.
\item[{Dok.~313}] Pascher, J., \emph{Kein Anfang: Zeitzyklus und
$\Omega_m^*$}. August 2026.
\item[{A040}] Pascher, J., \emph{Fraktale Korrektur}. A-Serie, 2026.
\item[{A050}] Pascher, J., \emph{Rekursion}. A-Serie, 2026.
\item[{A265}] Pascher, J., \emph{Rotverschiebung als statischer
Geometrie-Effekt}. A-Serie, 2026.
\item[{A271}] Pascher, J., \emph{Landauer und die Phasenraum-Buchhaltung}.
A-Serie, Juli 2026.
\item[{A280--A282}] Pascher, J., \emph{Phasenraum-Buchhaltung, Träger
und Information, Die Rechenkugel}. A-Serie, Juli 2026.
\end{description}

\section{Verwendung von KI-Werkzeugen}

Dieses Dokument ist unter Verwendung KI-gestützter Werkzeuge entstanden.
Fragestellungen, Setzungen und inhaltliche Entscheidungen stammen vom
Autor; die Werkzeuge formulieren aus und ordnen.
Verantwortung für Inhalt und Fehler liegt vollständig beim Autor.
Ausführlicher Wortlaut in A010.
